%!TEX root = TheoreticalPhysicsC.tex

\section{Lecture: The Harmonic Oscillator}
The energy of a \emph{classical} one-dimensional oscillator is given by
\begin{equation}
	E = \frac{1}{2m}p^2 + \frac{m\omega^2}{2}x^2,
\end{equation}
where $\omega$ is the frequency of the classical oscillator with the same
potential.

According to our quantisation table, the Hamiltonian of the quantum mechanical oscillator reads:
\begin{equation}
	\widehat{H} = \frac{1}{2m}\widehat{p}^2 + \frac{m\omega^2}{2}\widehat{x}^2.
\end{equation}	


The oscillator turns up in many contexts in physics: namely whenever one treats small motions about an equilibrium position approximately. A free quantum field can be analysed as a system of infinitely
many oscillators. It is worth getting to know this system thoroughly, and for that reason we shall look at it from several points of view.

Please do not picture a tiny little spring. As an \emph{effective} model of small displacements the oscillator is an approximation; for light quanta and for some optomechanical systems --- minute ``drums'' machined out of silicon --- it holds \emph{exactly}. Quantise the field of a resonator and you land precisely here: this system carries the whole of quantum optics.

The most direct method of obtaining the spectrum rests on the observation that $H$ can (essentially) be written as the square of a complex linear combination of $\widehat{p}$ and $\widehat{x}$. We set
\begin{equation}
	\widehat{a} = \sqrt{\frac{m\omega}{2\hbar}}\widehat{x} + i \sqrt{\frac{1}{2m\omega\hbar}}\widehat{p},
\end{equation}
so that
\begin{align}
	\widehat{a} \widehat{a}^\dag &= \frac{m\omega}{2\hbar}\widehat{x}^2 + \frac{1}{2m\omega\hbar}\widehat{p}^2 + \frac{i}{2\hbar} (\widehat{p}\widehat{x} - \widehat{x}\widehat{p}) \\
	\widehat{a}^\dag \widehat{a} &= \frac{m\omega}{2\hbar}\widehat{x}^2 + \frac{1}{2m\omega\hbar}\widehat{p}^2 - \frac{i}{2\hbar} (\widehat{p}\widehat{x} - \widehat{x}\widehat{p}).
\end{align}
Two warnings. First, it is by no means obvious that one should write down such a thing: the definition of $\widehat{a}$ comes out of nowhere and justifies itself solely through the identities that follow in a moment. Second, $\widehat{a}$ is not self-adjoint, and so is \emph{not an observable} --- no detector measures it; it is merely an aid to diagonalisation.

Using the canonical commutation relations $[\widehat{p},\widehat{x}] = -i\hbar\mathbb{I}$ we thus have
\begin{equation}\label{eq:ccra}
	\eqbox{ [\widehat{a}, \widehat{a}^\dag] = \mathbb{I}}
\end{equation}
and
\begin{equation}\label{eq:qho}
	\eqbox{\widehat{H} = \hbar \omega \left(\widehat{a}^\dag \widehat{a} + \tfrac{1}{2}\mathbb{I}\right).}
\end{equation}
The gain lies in the shape: an operator is positive exactly when it can be written as $\widehat{M}^\dag\widehat{M}$ --- and that is precisely what stands there explicitly in (\ref{eq:qho}), whereas in the original form it stays hidden.

To begin with, it is clear from this form of $H$ that $H \ge \hbar\omega/2 \mathbb{I}$. The oscillator is
therefore never ``at rest'': because of the uncertainty relation, $x=0$ and $p = 0$ cannot both be prepared to arbitrary accuracy. 

Is the ``zero-point energy'' $\hbar \omega/2$ itself an eigenvalue? The corresponding eigenvector $|\phi_0\rangle$ would have to satisfy the equation 
\begin{equation}
	\widehat{a}^\dag \widehat{a}|\phi_0\rangle = 0,
\end{equation}
that is, $\|\widehat{a}|\phi_0\rangle\|=0$, and hence $\widehat{a}|\phi_0\rangle = 0$. In the position representation this
is the ordinary differential equation $\langle x|\widehat{a}|\phi_0\rangle = 0$:
\begin{equation}
	\sqrt{\frac{m\omega}{2\hbar}}x\phi_0(x) + i\sqrt{\frac{1}{2m\omega\hbar}} \frac{\hbar}{i}\frac{d}{dx} \phi_0(x) = 0,
\end{equation}
that is, 
\begin{equation}
	\frac{d\phi_0}{dx}  = -\frac{m\omega}{\hbar} x\phi_0,
\end{equation}
with the solution
\begin{equation}
	\phi_0(x) = Ne^{-\frac{m\omega}{\hbar} \frac{x^2}{2}},
\end{equation}
where $N$ is the normalisation 
\begin{equation}
	N = \left(\frac{m\omega}{\pi\hbar}\right)^{\frac{1}{4}}
\end{equation}

Now one sees what all the effort was good for: by way of $\widehat{a}|\phi_0\rangle = 0$, a differential equation of \emph{second} order has become one of \emph{first} order. Until it is normalised, $\phi_0$ is only a function; it becomes a \emph{state} only through $\|\phi_0\| = 1$.

In deriving the uncertainty relation we used, exactly as here, the \emph{positivity} of the operator $\widehat{a}^\dag \widehat{a}$. There too the case of equality leads to the condition $\widehat{a}|\Psi\rangle = 0$. The ground state of the oscillator is therefore precisely a state of minimal uncertainty.

\begin{greybox}[Product rule for commutators]
In the next step a rule comes into play that you may not yet know. For arbitrary operators $\widehat{A}, \widehat{B}, \widehat{C}$ one has
\begin{equation}\label{eq:leibnizkomm}
	[\widehat{A}\widehat{B},\widehat{C}] = \widehat{A}[\widehat{B},\widehat{C}] + [\widehat{A},\widehat{C}]\widehat{B}
	\qquad \text{and} \qquad
	[\widehat{A},\widehat{B}\widehat{C}] = [\widehat{A},\widehat{B}]\widehat{C} + \widehat{B}[\widehat{A},\widehat{C}].
\end{equation}
The first identity is checked in a single line by adding and subtracting $\widehat{A}\widehat{C}\widehat{B}$:
\begin{equation}
	[\widehat{A}\widehat{B},\widehat{C}] = \widehat{A}\widehat{B}\widehat{C} - \widehat{C}\widehat{A}\widehat{B}
	= \widehat{A}\widehat{B}\widehat{C} - \widehat{A}\widehat{C}\widehat{B} + \widehat{A}\widehat{C}\widehat{B} - \widehat{C}\widehat{A}\widehat{B}
	= \widehat{A}[\widehat{B},\widehat{C}] + [\widehat{A},\widehat{C}]\widehat{B}.
\end{equation}
The proof of the second we leave to you. The commutator thus acts like a derivative --- only one has to keep an eye on the ordering.
\end{greybox}

For the calculation that follows we need the relations
\begin{align}
	[\widehat{H},\widehat{a}^\dag] &= \hbar \omega \widehat{a}^\dag \\
	[\widehat{H},\widehat{a}] &= -\hbar \omega \widehat{a}, 
\end{align}
which follow from (\ref{eq:ccra}) and (\ref{eq:qho}). The summand $\tfrac{1}{2}\hbar\omega\mathbb{I}$ in $\widehat{H}$ commutes with everything and drops out; with (\ref{eq:leibnizkomm}) there remains
\begin{equation}
	[\widehat{H},\widehat{a}^\dag] = \hbar\omega[\widehat{a}^\dag\widehat{a},\widehat{a}^\dag]
	= \hbar\omega\left(\widehat{a}^\dag[\widehat{a},\widehat{a}^\dag] + [\widehat{a}^\dag,\widehat{a}^\dag]\widehat{a}\right)
	= \hbar\omega \widehat{a}^\dag,
\end{equation}
and the second relation goes in just the same way.

If now $\widehat{H}|\Psi\rangle = E|\Psi\rangle$ is an eigenvector, then 
\begin{equation}
	\widehat{H}\widehat{a}^\dag|\Psi\rangle = (E  + \hbar \omega)\widehat{a}^\dag|\Psi\rangle
\end{equation}
So $\widehat{a}^\dag|\Psi\rangle$ satisfies the eigenvalue equation with eigenvalue $E + \hbar \omega$. One checks this by adding and subtracting $\widehat{a}^\dag\widehat{H}|\Psi\rangle$ --- a trick used again and again, for in this way a familiar commutator appears. Because 
\begin{equation}
	\|\widehat{a}^\dag|\Psi\rangle\|^2 = \langle \Psi|\widehat{a}\widehat{a}^\dag|\Psi\rangle = \langle \Psi|\widehat{a}^\dag\widehat{a}+\mathbb{I}|\Psi\rangle \ge \langle \Psi|\Psi\rangle \ge 0
\end{equation}
the vector $\widehat{a}^\dag|\Psi\rangle$ cannot be zero either: we really do
obtain a new eigenvector. Starting from $|\phi_0\rangle$ we obtain in this way the
eigenvectors 
\begin{equation}\label{eq:qhoeigenstates}
	\eqbox{|\phi_n\rangle = \frac{(\widehat{a}^\dag)^n|\phi_0\rangle}{\|(\widehat{a}^\dag)^n|\phi_0\rangle \|}}
\end{equation}
and eigenvalues
\begin{equation}
	\eqbox{E_n = \hbar\omega(n+\tfrac{1}{2}).}
\end{equation}

It is because of this procedure, incidentally, that the operators $\widehat{a}$ and $\widehat{a}^\dag$ are also called \emph{ladder operators}, or ``lowering operator'' and ``raising operator'' respectively. We use the lowering operator $\widehat{a}$ to show that there are no further eigenvalues either. For if $|\Psi\rangle$ is an eigenvector with eigenvalue $E$, then $(\widehat{a})^n|\Psi\rangle$ satisfies the eigenvalue equation with eigenvalue $E - n\hbar \omega$. For large enough $n$ this expression becomes negative, and so cannot be an eigenvalue of the positive operator
$H$. For large $n$ we must therefore have $(\widehat{a})^n|\Psi\rangle=0$. Let $n$ be the smallest number for which this holds. Then $|\Phi\rangle = (\widehat{a})^{n-1}|\Psi\rangle$ is a non-zero solution
of the equation $\widehat{a}|\Phi\rangle = 0$, and hence, by our calculation above, proportional to $|\phi_0\rangle$. By the commutation relation for $\widehat{a}$ and $\widehat{a}^\dag$, $|\Psi\rangle$ is then a multiple of
$(\widehat{a}^\dag)^n|\Phi_0\rangle$. The complete list of eigenfunctions given here is therefore a basis of the Hilbert space.

The denominator in (\ref{eq:qhoeigenstates}) is not a blemish: $\widehat{a}^\dag|\Psi\rangle$ is indeed an eigenvector again, but it is not normalised; only the division by the norm turns it into a state. Do not expect, by the way, that the procedure carries over in general --- essentially it works only here. The ladder operators made their real career under different names, as \emph{creation} and \emph{annihilation operators} in many-body theory.

The norm $\|(\widehat{a}^\dag)^n|\Phi_0\rangle\|$ can likewise be computed from the commutation relation
alone. We have
\begin{equation}
	\|\widehat{a}^\dag|\phi_n\rangle\|^2 = \langle\phi_n|\widehat{a}\widehat{a}^\dag|\phi_n\rangle =  \langle\phi_n|\left(\frac{1}{\hbar\omega}\widehat{H}+\frac12\mathbb{I}\right)|\phi_n\rangle=(n+1)\|\phi_n\|^2 = (n+1).
\end{equation}
Hence 
\begin{equation}
	\widehat{a}^\dag|\phi_n\rangle = \sqrt{n+1}|\phi_{n+1}\rangle,
\end{equation}
and (by induction)
\begin{equation}
	|\phi_n\rangle = \frac{1}{\sqrt{n!}}(\widehat{a}^\dag)^n|\Phi_0\rangle.
\end{equation}

Since the $|\phi_n\rangle$ form a basis of the Hilbert space, every operator can be expressed in terms of its matrix elements with respect to this basis. This is worth doing whenever time-dependent quantities are to be studied, since the time evolution operators are then multiplication operators. We obtain
(with the matrix notation $[\widehat{C}]_{nm} = \langle\phi_n|\widehat{C}|\phi_m\rangle$)
\begin{equation}
	[\widehat{a}^\dag]_{nm} = \sqrt{n}\delta_{n,m+1} \quad \text{and} \quad [\widehat{a}]_{nm} = \sqrt{m}\delta_{n+1,m}.
\end{equation}
With 
\begin{equation}
	\widehat{p} = \frac{\sqrt{2m\omega\hbar}}{2i}(\widehat{a}-\widehat{a}^\dag) 
\end{equation}
one obtains at once the matrix for $\widehat{p}$, and correspondingly the matrix for $\widehat{x}$. Written out in full, they are:
\begin{equation}
	\widehat{p} = \frac{\sqrt{2m\omega\hbar}}{2i}
	\begin{pmatrix}
		0 & \sqrt{1} & 0 & 0 & 0 & 0 & \hdots \\
		-\sqrt{1} & 0 & \sqrt{2} & 0 & 0 & 0 & \\
		0 & -\sqrt{2} & 0 & \sqrt{3} & 0 & 0 & \\
		0 & 0 & -\sqrt{3} & 0 & \sqrt{4} & 0 & \\
		0 & 0 & 0 & -\sqrt{4} & 0 & \sqrt{5} & \\
		0 & 0 & 0 & 0 & -\sqrt{5} & 0 & \\
		\vdots & & & & & & \ddots \\
	\end{pmatrix}
\end{equation}
and
\begin{equation}
	\widehat{x} = \sqrt{\frac{\hbar}{2m\omega}}
	\begin{pmatrix}
		0 & \sqrt{1} & 0 & 0 & 0 & 0 & \hdots \\
		\sqrt{1} & 0 & \sqrt{2} & 0 & 0 & 0 & \\
		0 & \sqrt{2} & 0 & \sqrt{3} & 0 & 0 & \\
		0 & 0 & \sqrt{3} & 0 & \sqrt{4} & 0 & \\
		0 & 0 & 0 & \sqrt{4} & 0 & \sqrt{5} & \\
		0 & 0 & 0 & 0 & \sqrt{5} & 0 & \\
		\vdots & & & & & & \ddots \\
	\end{pmatrix}.
\end{equation}

An important advantage of the \emph{algebraic method} of the ladder operators mentioned above is
that the results are immediately available in every representation of the Hilbert space. For example, formula (\ref{eq:qhoeigenstates}) delivers the wave functions/eigenfunctions $\phi_n(x) \in L^2(\mathbb{R})$ in the position representation as soon as $\widehat{a}$ is written out explicitly in that representation. By induction one sees that 
\begin{equation}
	\phi_n(x) = c_n \mathrm{H}_n(\lambda x)  e^{-\frac{\lambda^2 x^2}{2}}
\end{equation}
with
\begin{equation}
	\lambda = \sqrt{\frac{m\omega}{\hbar}}\quad \text{and} \quad c_n = \sqrt{\frac{\lambda}{2^n n!\sqrt{\pi}}},
\end{equation}
where
$\mathrm{H}_n(\lambda x)$ are the so-called \emph{Hermite polynomials}. The first $3$ Hermite polynomials
read
\begin{align}
	\mathrm{H}_0(z) &= 1,\\
	\mathrm{H}_1(z) &= 2z,\\
	\mathrm{H}_2(z) &= 4z^2-2,
\end{align}
and so on.

% fig_vl12_leiter.tex — Der harmonische Oszillator: Eigenfunktionen in der
% Parabel und dieselben Niveaus als Leiter mit a und a^dagger.
%
% Follows the conventions of fig_vl11_kasten (the reference plate); the two
% figures are meant to be read against each other, box versus oscillator.
%
%   * geometry and curves come from the closed-form solution derived in the
%     lecture, never from a freehand sketch.  In the natural units
%     xi = x sqrt(m omega / hbar) and E/(hbar omega):
%         V(xi)     = xi^2 / 2                        (the parabola)
%         E_n       = n + 1/2                         (equally spaced)
%         phi_n(xi) = c_n H_n(xi) exp(-xi^2/2),
%                     c_n = (2^n n! sqrt(pi))^{-1/2}
%                     c_0..c_3 = 0.751126, 0.531126, 0.265562, 0.108416
%     Consequences that are therefore forced rather than drawn by eye:
%     the node count (phi_n has n zeros), the parity, the constant gap,
%     and the classical turning points E_n = V(xi)  =>  xi = ±sqrt(2n+1),
%     which is why the grey level rules get longer like sqrt(2n+1) and
%     end exactly on the parabola.  The plot corners are chosen so that
%     the parabola exits precisely there: ymax = xmax^2/2 = 5.12.
%   * amplitude factor 0.6: the largest |phi_n| is phi_0(0) = 0.7511, so
%     0.6 keeps every curve within 0.45 of its own level, i.e. 45% of the
%     gap hbar omega — the same clearance the box figure uses (1.3 of 3).
%     No curve touches a neighbouring level line.
%   * colour encodes the quantum number n (the ordered ramp of `tpclevels`),
%     in both panels and nothing else: rung n in (b) has the colour of
%     curve n in (a).  Entries 1, 2, 4, 5 (navy, teal, orange, red);
%     entry 3, munsell, is skipped on purpose — beside cgblue it is the
%     same teal to the eye, so it would break the very code it is
%     supposed to carry.  This is the choice fig_vl11_kasten makes, and
%     the two plates are meant to be read against each other.  Level
%     rules in (a) are neutral grey — reference lines, not data — and so
%     are the turning-point dots.
%   * both axes share ymin/ymax and height, so E_n in (a) sits at exactly
%     the height of |phi_n> in (b): the panels are one energy axis, read
%     twice.
%   * the classically forbidden region E < V is the area *under* the
%     parabola; it is hatched, not filled, so it survives greyscale
%     printing, and one sees the tails of phi_n reaching into it — the
%     one thing the infinite well cannot show.
%
% Panel (a) phi_0..phi_3 on their levels inside the parabola.
% Panel (b) the same levels as a ladder: a^dagger up, a down, both with
%           weight sqrt(n+1) on the gap between phi_n and phi_{n+1}, and
%           the bottom rung where a|phi_0> = 0 ends the ladder.

\begin{center}
\begin{tikzpicture}

% ─────────────────────────────────────────────────────────────
%  (a)  Eigenfunctions on their levels inside the parabola
% ─────────────────────────────────────────────────────────────
\begin{axis}[
    name=parabel,
    tpclevelplot,
    width=7.2cm, height=5.8cm,
    xmin=-3.2, xmax=3.2,
    ymin=0,    ymax=5.12,
    title={(a) \tpcifenglish{$\phi_n$ in the parabola}%
                            {$\phi_n$ in der Parabel}},
    domain=-3.2:3.2, samples=201,
]

  % Classically forbidden region E < V: everything below the parabola.
  \addplot[forbidden, draw=none, samples=161] {x^2/2} \closedcycle;

  % The potential itself.  Neutral heavy grey, not the usual dark ink:
  % the parabola is the boundary of the hatched region, i.e. reference
  % structure, and spacecadet is already taken by the n=0 eigenfunction —
  % two navy curves crossing at the turning point cannot be told apart.
  \addplot[potential, birrengray!85!black, line width=1.4pt] {x^2/2};

  % Energy levels E_n = (n + 1/2) hbar omega, drawn from turning point to
  % turning point, xi = ±sqrt(2n+1): neutral reference rules.
  \addplot[levellinegrey, domain=-1.0000:1.0000, samples=2] {0.5};
  \addplot[levellinegrey, domain=-1.7321:1.7321, samples=2] {1.5};
  \addplot[levellinegrey, domain=-2.2361:2.2361, samples=2] {2.5};
  \addplot[levellinegrey, domain=-2.6458:2.6458, samples=2] {3.5};

  % Classical turning points, where the levels meet the parabola
  \addplot[only marks, mark=*, mark size=1.2pt, birrengray, draw=none]
    coordinates {(-1.0000,0.5) (1.0000,0.5) (-1.7321,1.5) (1.7321,1.5)
                 (-2.2361,2.5) (2.2361,2.5) (-2.6458,3.5) (2.6458,3.5)};

  % Eigenfunctions about their own levels.  Each is drawn out towards one
  % unit beyond its own turning point, where the tail has decayed to a few
  % per cent — and, for the upper two, as far beyond it as the plate
  % allows while still ending clear of the boundary.  The curves then nest
  % like the level rules they belong to, no state trails a long flat line
  % across the whole plate, and none runs off the edge.
  \addplot[wavefn, spacecadet,   domain=-2.10:2.10]
    {0.5 + 0.6*0.751126*exp(-x^2/2)};
  \addplot[wavefn, cgblue,       domain=-2.60:2.60]
    {1.5 + 0.6*0.531126*(2*x)*exp(-x^2/2)};
  \addplot[wavefn, birrenorange, domain=-2.85:2.85]
    {2.5 + 0.6*0.265562*(4*x^2-2)*exp(-x^2/2)};
  \addplot[wavefn, birrenred,    domain=-3.05:3.05]
    {3.5 + 0.6*0.108416*x*(8*x^2-12)*exp(-x^2/2)};

  % Level labels, in the left margin.  Written out rather than looped:
  % \foreach defers expansion inside a pgfplots axis.
  \node[lbl, anchor=east] at (axis cs:-3.3,0.5) {$E_{0}$};
  \node[lbl, anchor=east] at (axis cs:-3.3,1.5) {$E_{1}$};
  \node[lbl, anchor=east] at (axis cs:-3.3,2.5) {$E_{2}$};
  \node[lbl, anchor=east] at (axis cs:-3.3,3.5) {$E_{3}$};

  % Node counts: phi_n has n zeros
  \node[lbl, text=birrengray, font=\scriptsize, anchor=west]
    at (axis cs:3.3,0.5) {$n=0$};
  \node[lbl, text=birrengray, font=\scriptsize, anchor=west]
    at (axis cs:3.3,1.5) {$n=1$};
  \node[lbl, text=birrengray, font=\scriptsize, anchor=west]
    at (axis cs:3.3,2.5) {$n=2$};
  \node[lbl, text=birrengray, font=\scriptsize, anchor=west]
    at (axis cs:3.3,3.5) {$n=3$};

  % Axis furniture
  % The potential is named in its own colour, the hatched wedge in the
  % grey of the hatching: label and object match, so neither needs a leader.
  \node[lbl, text=birrengray!85!black] at (axis cs:0,4.65)
    {$V(x)=\tfrac{1}{2}m\omega^2x^2$};
  % Named on clean paper: grey glyphs on grey hatching are unreadable, so
  % the node punches a small white hole for itself.
  \node[lbl, text=birrengray, font=\scriptsize, fill=white, inner sep=1.6pt]
    at (axis cs:-2.30,0.28) {$E<V$};
  \node[lbl, anchor=north] at (axis cs:0,-0.24) {$x\sqrt{m\omega/\hbar}$};
\end{axis}

% ─────────────────────────────────────────────────────────────
%  (b)  The same levels as a ladder
% ─────────────────────────────────────────────────────────────
\begin{axis}[
    tpclevelplot,
    at={(parabel.south east)}, anchor=south west, xshift=16mm,
    width=4.6cm, height=5.8cm,
    xmin=-1.0, xmax=1.75,
    ymin=0,    ymax=5.12,
    title={(b) \tpcifenglish{ladder operators}{Leiteroperatoren}},
]

  % E = 0 and the rungs.  Colour = n, exactly as in panel (a).
  \draw[classicallim] (axis cs:-0.62,0) -- (axis cs:1.25,0);
  \draw[wavefn, spacecadet]   (axis cs:0,0.5) -- (axis cs:1.25,0.5);
  \draw[wavefn, cgblue]       (axis cs:0,1.5) -- (axis cs:1.25,1.5);
  \draw[wavefn, birrenorange] (axis cs:0,2.5) -- (axis cs:1.25,2.5);
  \draw[wavefn, birrenred]    (axis cs:0,3.5) -- (axis cs:1.25,3.5);
  \draw[levellinegrey]        (axis cs:0,4.5) -- (axis cs:1.25,4.5);
  \node[lbl, text=birrengray, font=\large] at (axis cs:0.625,4.82) {$\vdots$};

  % a^dagger raises, a lowers; the weight sqrt(n+1) belongs to the gap.
  \draw[-{Stealth[length=5pt,width=3.5pt]}, spacecadet, line width=0.9pt]
    (axis cs:0.30,0.56) -- (axis cs:0.30,1.44);
  \draw[-{Stealth[length=5pt,width=3.5pt]}, spacecadet, line width=0.9pt]
    (axis cs:0.30,1.56) -- (axis cs:0.30,2.44);
  \draw[-{Stealth[length=5pt,width=3.5pt]}, spacecadet, line width=0.9pt]
    (axis cs:0.30,2.56) -- (axis cs:0.30,3.44);
  \draw[-{Stealth[length=5pt,width=3.5pt]}, spacecadet, line width=0.9pt]
    (axis cs:0.30,3.56) -- (axis cs:0.30,4.44);

  \draw[-{Stealth[length=5pt,width=3.5pt]}, spacecadet, line width=0.9pt]
    (axis cs:0.95,1.44) -- (axis cs:0.95,0.56);
  \draw[-{Stealth[length=5pt,width=3.5pt]}, spacecadet, line width=0.9pt]
    (axis cs:0.95,2.44) -- (axis cs:0.95,1.56);
  \draw[-{Stealth[length=5pt,width=3.5pt]}, spacecadet, line width=0.9pt]
    (axis cs:0.95,3.44) -- (axis cs:0.95,2.56);
  \draw[-{Stealth[length=5pt,width=3.5pt]}, spacecadet, line width=0.9pt]
    (axis cs:0.95,4.44) -- (axis cs:0.95,3.56);

  \node[lbl] at (axis cs:0.625,1.0) {$\sqrt{1}$};
  \node[lbl] at (axis cs:0.625,2.0) {$\sqrt{2}$};
  \node[lbl] at (axis cs:0.625,3.0) {$\sqrt{3}$};
  \node[lbl] at (axis cs:0.625,4.0) {$\sqrt{4}$};

  \node[lbl, anchor=east] at (axis cs:0.24,2.0) {$\widehat{a}^\dag$};
  \node[lbl, anchor=west] at (axis cs:1.01,2.0) {$\widehat{a}$};

  % The gap and the zero-point offset, measured off to the left
  \draw[classicallim] (axis cs:0,0.5) -- (axis cs:-0.62,0.5);
  \draw[classicallim] (axis cs:0,1.5) -- (axis cs:-0.62,1.5);
  \draw[{Stealth[length=4pt]}-{Stealth[length=4pt]}, spacecadet,
        line width=0.6pt] (axis cs:-0.45,0.5) -- (axis cs:-0.45,1.5);
  \draw[{Stealth[length=4pt]}-{Stealth[length=4pt]}, spacecadet,
        line width=0.6pt] (axis cs:-0.45,0.02) -- (axis cs:-0.45,0.48);
  \node[lbl, anchor=east] at (axis cs:-0.50,1.0)  {$\hbar\omega$};
  \node[lbl, anchor=east] at (axis cs:-0.50,0.25) {$\tfrac{1}{2}\hbar\omega$};

  % State labels
  \node[lbl, anchor=west] at (axis cs:1.31,0.5) {$|\phi_0\rangle$};
  \node[lbl, anchor=west] at (axis cs:1.31,1.5) {$|\phi_1\rangle$};
  \node[lbl, anchor=west] at (axis cs:1.31,2.5) {$|\phi_2\rangle$};
  \node[lbl, anchor=west] at (axis cs:1.31,3.5) {$|\phi_3\rangle$};
  \node[lbl, text=birrengray, font=\scriptsize, anchor=west]
    at (axis cs:1.31,0) {$E=0$};

  % Legend for the arrows: the rule the weights come from.
  \node[lbl, font=\scriptsize, align=center, anchor=north]
    at (axis cs:0.35,-0.20)
    {$\widehat{a}^\dag|\phi_n\rangle=\sqrt{n+1}\,|\phi_{n+1}\rangle$\\[-1pt]
     $\widehat{a}|\phi_0\rangle=0$};
\end{axis}

\end{tikzpicture}
\end{center}


On the left, each grey level line ends where the parabola cuts it: at the
classical turning point, beyond which $\phi_n$ nevertheless still reaches into
the hatched region. On the right the same levels as a ladder ---
equally spaced, and breaking off at the bottom.

Before we come to the time evolution, a warning.

\begin{merksatz}[Take care when exponentiating operators]
For operators, $e^{\widehat{A}+\widehat{B}} = e^{\widehat{A}}e^{\widehat{B}}$ holds \emph{only} when $[\widehat{A},\widehat{B}] = 0$.
\end{merksatz}

It is easy to see where this fails: already the quadratic term no longer agrees, since
\begin{equation}
	(\widehat{A}+\widehat{B})^2 = \widehat{A}^2 + \widehat{A}\widehat{B} + \widehat{B}\widehat{A} + \widehat{B}^2
	\neq \widehat{A}^2 + 2\widehat{A}\widehat{B} + \widehat{B}^2.
\end{equation}
If, on the other hand, the two commute, then they behave like numbers and the series agree term by term. For the oscillator we are allowed to do so: with $\widehat{P}_n = |\phi_n\rangle\langle\phi_n|$ we have $\widehat{H} = \sum_n E_n\widehat{P}_n$, and for $n \neq m$
\begin{equation}
	[\widehat{P}_n,\widehat{P}_m] = |\phi_n\rangle\langle\phi_n|\phi_m\rangle\langle\phi_m|
	- |\phi_m\rangle\langle\phi_m|\phi_n\rangle\langle\phi_n| = 0
\end{equation}
because of the orthonormality $\langle\phi_n|\phi_m\rangle = \delta_{nm}$. That is exactly what justifies writing $U_t$ down term by term.

With the diagonalisation of the Hamiltonian we have of course also gained detailed information about the time evolution of the oscillator. Like the classical oscillator, the quantum mechanical one is periodic with period $T = 2\pi/\omega$. This follows from the list of eigenvalues alone: by the functional calculus the eigenvalues of the time evolution operator $U_T = e^{-\frac{iT}{\hbar} \widehat{H}}$ are
\begin{equation}
	e^{-\frac{iT}{\hbar} \hbar\omega(n+\frac12)} = e^{-2\pi i (n+\frac12)} = -1,
\end{equation}
so independent of $n$, and $U_T = -\mathbb{I}$ is a multiple of the identity operator $\mathbb{I}$.

Thus for the time dependence of all expectation values, that is, by Postulate 5, for the functions of the form $f(t) = \tr(U_t \rho U_t^\dag A)$:
\begin{equation}
	f(t+T) = \tr(U_tU_T \rho U_T^\dag U_t^\dag A) =f(t).
\end{equation}	
The minus sign in $U_T = -\mathbb{I}$ does no harm: in $\rho_T = U_T\rho_0U_T^\dag$ the signs cancel, the global phase drops out of its own accord --- an advantage of the density matrix formalism. Nothing at all was assumed about the initial state.

The analogy with the classical evolution goes much further still. The simplest way to see this here is to compute the time dependence of the ladder operator $\widehat{a}$. We have 
\begin{equation}
	\frac{d}{dt} U_t^\dag \widehat{a} U_t = U_t^\dag \left(\frac{i}{\hbar} [\widehat{H},\widehat{a}] \right) U_t = -i\omega U_t^\dag \widehat{a} U_t
\end{equation}
hence
\begin{equation}
	U_t^\dag \widehat{a} U_t = e^{-i\omega t}\widehat{a}.
\end{equation}
Along the way one finds that the quantum mechanical position and momentum operators follow exactly the classical equations of motion:
\begin{equation}
	\begin{split}
		\widehat{p}(t) &= U_t^\dag \widehat{p} U_t  = \cos(\omega t) \widehat{p}  - m\omega\sin(\omega t) \widehat{x} \\
		\widehat{x}(t) &= U_t^\dag \widehat{x} U_t  = \frac{1}{m\omega}\sin(\omega t) \widehat{p}  + \cos(\omega t) \widehat{x}
	\end{split}
\end{equation}

We are working here in the \emph{Heisenberg picture}: it is not the states that move but the measurements, and $U_t^\dag \widehat{A} U_t$ is ``the same measurement, $t$ seconds later''. With that the oscillator is finished --- there are whole books about this one system.

